\documentclass[11pt]{article}
\usepackage{amsmath, amssymb}
\usepackage{geometry}
\usepackage{array}
\usepackage{booktabs}
\geometry{margin=0.7in}
\setlength{\parindent}{0pt}
\setlength{\parskip}{0.35em}
\renewcommand{\arraystretch}{1.15}

\newcommand{\cheatsheettitle}[2]{%
\begin{center}
{\LARGE #1}\\[0.4em]
{\large #2}
\end{center}
}


\begin{document}
\cheatsheettitle{Algebra I Proportional Relationships Cheatsheet}{Module 4}

\section*{Ratios and Rates}
A ratio compares two quantities.
\[
a:b,\qquad \frac{a}{b},\qquad a\text{ to }b
\]
A unit rate has denominator $1$.
\[
\frac{\text{miles}}{\text{hours}}=\text{miles per hour}
\]

\section*{Proportions}
\[
\frac{a}{b}=\frac{c}{d}\qquad (b\ne0,\ d\ne0)
\]
Cross products are equal:
\[
ad=bc
\]

\section*{Direct Variation}
\[
y=kx
\]
The constant of proportionality is
\[
k=\frac{y}{x}
\]
Graph feature:
\[
y=kx\text{ passes through }(0,0)
\]

\section*{Slope as Rate of Change}
\[
m=\frac{\text{change in }y}{\text{change in }x}
=\frac{y_2-y_1}{x_2-x_1}
\]
For proportional relationships:
\[
m=k
\]

\section*{Percent Change}
\[
\text{percent change}=\frac{\text{new}-\text{old}}{\text{old}}\cdot100\%
\]
Increase:
\[
\text{new}=\text{old}(1+r)
\]
Decrease:
\[
\text{new}=\text{old}(1-r)
\]
where $r$ is written as a decimal.

\section*{Tables and Graphs}
A relationship is proportional when:
\begin{itemize}
  \item the ratio $\frac{y}{x}$ is constant for all nonzero $x$,
  \item the graph is a line,
  \item the line passes through the origin.
\end{itemize}

\section*{Quick Checks}
\begin{itemize}
  \item Keep units attached to rates.
  \item Proportional means linear through the origin.
  \item A constant difference is linear; a constant ratio through the origin is proportional.
  \item Convert percents to decimals before multiplying.
\end{itemize}

\end{document}
